% =============================================================
%  Evaluation.tex —— §11 模型评价
%  数据来源：四问素材库 + 实验日志 + 真机汇总
%  编译方式：xelatex Evaluation.tex
% =============================================================
\documentclass[UTF8,12pt,a4paper]{ctexart}

\usepackage{amsmath}
\usepackage{amssymb}
\usepackage{geometry}
\geometry{margin=2.5cm}
\usepackage{booktabs}
\usepackage{float}
\usepackage{caption}
\usepackage{enumitem}

\setcounter{section}{10}

\begin{document}

% =============================================================
\section{模型评价}
% =============================================================

% -------------------------------------------------------------
\subsection{模型优势}
% -------------------------------------------------------------

\textbf{(1) 题目契合度高。} 四问均围绕 QUBO 建模、SDK/CIM 求解与量子比特最少化展开：Q1/Q2 采用 $n^2=225$ 比特 one-hot 编码完整覆盖小规模 TSP 与 TSPTW；Q3 通过滚动窗口分解使 $n=50$ 大规模问题的子 QUBO 维度 $\le 289$ 比特；Q4 多车独立子 QUBO 维度 $\le 64$ 比特。所有子 QUBO 均严格满足比特预算原则与真机规模约束（$\le 550$ 比特），可直接上传 CPQC-550 真机求解。

\textbf{(2) 模型递进清楚。} 全文按 ``TSP $\to$ TSPTW $\to$ 大规模分解 $\to$ 多车容量约束'' 的顺序逐步引入难点：Q1 解决路径基线，Q2 引入软时间窗惩罚，Q3 解决变量爆炸下的可解性问题，Q4 进一步引入车辆数与载重约束。每一问相对前一问明确给出新增的决策变量、约束与目标项，模型层次分明。

\textbf{(3) 比特预算清晰。} 各问 QUBO 变量数均经过显式公式推导并与等价编码进行比较：Q1/Q2 为 $n^2=225$；Q3 单段子 QUBO 在 $n_{\mathrm{seg}}\le 17$ 时为 $\le 289$；Q4 单车子 QUBO 在 $n_k\le 8$ 时为 $\le 64$。全部子问题均低于 550 比特上限，并在比特数 vs 规模 $n$ 的曲线中得到可视化验证。

\textbf{(4) SDK--CIM 链路一致性强。} 通过统一的 \texttt{solve\_subqubo(Q, backend)} 接口，将 SDK 模拟退火验证与 CIM 真机调用封装在同一求解流程中：编码方式、比特预算与 8-bit 精度调整完全一致，仅求解器后端切换。Python / Kaiwu SDK SA / CIM 真机三方在 Q2（$J=84121$）、Q3（$J=4{,}941{,}906$，seg3 严格相等）等问题上得到相互印证，构成重要验证证据，增强了结果可信度。

\textbf{(5) 可复现性强。} 全文采用统一的 \texttt{load\_data}、\texttt{evaluate}、QUBO 构造与真机调用链路：所有结果可由 \texttt{src/01--15} 系列脚本在 Python 3.10 与 Kaiwu SDK 1.3.1 环境下复算；JSON 结果与图表落盘路径在论文素材库和实验日志中逐一登记，提高了模型可复现性。

% -------------------------------------------------------------
\subsection{模型不足}
% -------------------------------------------------------------

\textbf{(1) 全局最优性证明不足。} Q1 ($n=15$ TSP) 已用 Held--Karp 动态规划得到 $J=29$ 的精确验证；Q2--Q4 主要依赖多算法一致性、参数扫描与消融检验，尚不能构成严格的全局最优性证明，故按 ``近似最优解'' 表述。

\textbf{(2) SDK SA 原始可行率较低。} Q1$\sim$Q4 子 QUBO 在 SDK 模拟退火上原始可行率落在 $0.5\%\sim 2\%$ 区间。one-hot 强约束导致普通单翻转邻域较难直接采到可行解，因此本文依赖 ``采样--修复'' 流程托底；后续可探索块翻转或约束保持邻域以提升原始可行率。

\textbf{(3) 真机配额限制。} 玻色 CPQC-550 平台每次任务占用 1 次配额，全四问完整真机验证至少需要 12 次配额。Q3 seg1 因账号配额异常未获得有效返回，本文未补造结果，而是在结果表中如实标注；其余子问题真机结果均完整登记。

\textbf{(4) 启发式搜索仍存在局部最优风险。} 在部分车辆数设置（如 Q4 $K=10$）下，启发式起点池可能未覆盖更优邻域，使综合目标略劣于参考解；其余 $K$ 取值下解优于参考或与之持平。这说明大规模 CVRPTW 在当前混合框架下仍存在局部最优风险，需要进一步丰富邻域算子或引入多次重启策略。

\textbf{(5) 时间窗惩罚项占主导。} 当客户数量增大而时间窗较紧时（如 Q3 数据集 $b_i\le 30$、$\sum$服务$+$travel$\approx 156$），目标函数中惩罚项占比远超运输时间，使模型更侧重减少违约程度而非单纯压缩运输时间。这一现象由数据集自身的结构性特征决定，需通过引入多车辆（如 Q4）从根本上缓解。

\textbf{(6) 客户分配未与车内排序联合优化。} Q4 客户分配 $y_{i,k}$ 由启发式 warm start 给定，未与车内排序统一进入 QUBO，因此本文模型仍不是完全联合优化模型；当数据集需求或时间窗大幅变动时，需重新运行启发式分配池。

% -------------------------------------------------------------
\subsection{改进方向与推广价值}
% -------------------------------------------------------------

\textbf{改进方向。} 针对前述不足，可从以下五个方向继续推进：

\begin{enumerate}[label=(\arabic*),leftmargin=*]
    \item 引入精确 ILP 或 CP-SAT 求解器，对小中规模算例（如 $n\le 25$ 的 TSPTW、$K\le 4$ 的 CVRPTW）提供更强的最优性验证基线，从而进一步缩小 ``近似最优'' 与 ``全局最优'' 之间的论证差距；
    \item 设计块翻转（block-flip）或约束保持邻域，使 SDK 模拟退火能在不破坏 one-hot 约束的前提下进行采样，从而提高原始可行率，降低对采样--修复流程的依赖；
    \item 引入自适应罚系数策略（如根据当前违反程度动态调整 $A$、$T_0$），减少人工调参成本，并提升模型在不同数据分布下的鲁棒性；
    \item 探索客户分配 $y_{i,k}$ 与车内排序 $x_{i,t,k}$ 的联合 QUBO 编码，将 Q4 由两阶段流程升级为完全联合优化模型，代价是比特数上升；可结合分解策略缓解规模压力；
    \item 增加 CIM 真机的多次重复实验，统计哈密顿量演化与最优解分布的随机性，对真机求解的稳定性进行系统评估。
\end{enumerate}

\textbf{推广价值。} 本文所建模型与混合求解框架具有以下推广前景：

\begin{enumerate}[label=(\arabic*),leftmargin=*]
    \item \textbf{物流场景推广}：可推广到多仓库、多车型、动态订单与实时调度等更复杂的智慧物流场景，仅需在 QUBO 构造层面扩展决策变量与约束项；
    \item \textbf{方法层面推广}：one-hot 编码、采样--修复流程与统一 \texttt{solve\_subqubo} 接口可迁移到其他 QUBO 型组合优化问题，如作业车间调度、图着色、最大独立集等；
    \item \textbf{真机适配推广}：分段 subQUBO 与 block-diag 合并思路适用于当前 CIM 真机比特数有限条件下的大规模问题求解，可作为量子--经典混合优化在工业级算例上的通用范式。
\end{enumerate}

\end{document}
